\chapter{作品概述}

\section{应用背景与研究意义}

\subsection{定向通信中的双方身份约束}

数字化业务中的消息传输，往往同时包含接收权限和发送来源两类要求。以机构设备管理为例，运维
中心向某一受控终端下发配置指令时，希望只有该终端能够读取；终端收到指令后，还要确认消息来自
自己认可的运维中心，而不是来自另一名合法用户或被替换的发送节点。类似要求也存在于机构内部的
定向文件投递、应急通知、跨部门业务消息和来源受限的数据分发中。

传统系统通常分别使用公钥加密和数字签名完成这两项任务。发送端先取得接收者公钥或证书并加密，
再使用自己的签名密钥认证消息；接收端依次验证证书、检查签名并解密。该方式能够实现功能，但
“指定谁接收”和“接收者认可谁发送”分散在不同对象和处理步骤中。应用协议必须额外规定证书
传递、身份字段、签名对象、加密顺序和异常处理，才能保证两项身份条件始终绑定在同一条消息上。

身份基匹配加密（Identity-Based Matchmaking Encryption，IB-ME）面向这种双方身份约束。
发送者在加密时写入接收身份，接收者在解密时写入自己认可的发送身份；只有两项条件同时成立，
应用消息才会恢复。图~\ref{fig:application-background}以机构运维场景说明密钥和消息的关系。可信
发钥机构核验业务身份并发放角色密钥，公开网络只传输密文，业务端直接使用机构账号或设备编号
表达通信对象。

\begin{figure}[H]
\centering
\begin{tikzpicture}[node distance=10mm and 10mm]
 \node[check,text width=36mm] (authority) {可信发钥机构\\核验身份并发放角色密钥};
 \node[flow,text width=37mm,below left=of authority] (sender) {运维中心$s$\\指定接收终端$r$};
 \node[flow,text width=34mm,right=of sender] (channel) {公开网络\\传输IB-ME密文};
 \node[check,text width=39mm,right=of channel] (receiver) {受控终端$r'$\\只认可运维中心$s'$};
 \node[check,text width=41mm,below=of channel] (result) {$r=r'$且$s=s'$：执行指令\\任一身份不匹配：拒绝};
 \draw[flowarrow] (authority)--node[left,font=\small]{发送密钥}(sender);
 \draw[flowarrow] (authority)--node[right,font=\small]{接收密钥}(receiver);
 \draw[flowarrow] (sender)--(channel);
 \draw[flowarrow] (channel)--(receiver);
 \draw[flowarrow] (receiver)--(result);
\end{tikzpicture}
\caption{机构定向通信中的身份与密钥关系}
\label{fig:application-background}
\end{figure}

这一密码功能并不限定于单一行业。凡是由统一机构管理成员或设备身份、消息需要定向送达且接收端
需要限制来源的系统，都可以把IB-ME作为消息保护组件。不同场景对在线延迟、密文长度和离线准备
成本的要求不同，因此还需要能够更换底层算法、比较多种实现并按使用方式选择组合。第八章将结合
本作品的三组代表方案进一步分析应用模式和部署条件。

\subsection{量子计算背景下的公钥密码迁移}

公钥加密和数字签名是上述通信机制的基础。当前广泛使用的RSA、Diffie--Hellman和椭圆曲线密码
分别依赖整数分解或离散对数问题。Shor算法表明，具备足够规模的通用量子计算机可以在多项式时间
内求解这两类问题\cite{shor1994}。对于需要长期保护的数据和生命周期较长的信息系统，密码迁移
不仅要准备新的底层算法，还要提前解决身份管理、协议组合、对象格式、运行资源和软件接入问题。

格密码以高维格上的计算问题为基础，目前没有已知的经典或量子多项式时间算法能够求解其适当参数
实例。格上误差学习（LWE）、短整数解（SIS）及其环、模和NTRU结构变体，能够支持公钥加密、
数字签名、身份基加密和身份基签名。不同结构在对象尺寸、计算方式和实现复杂度上各有特点，因而
适合构成可比较、可替换的后量子密码组件。

身份基密码还为统一管理的机构系统提供了直接的身份表达。它允许邮箱、机构账号、设备编号等规范
身份参与密码运算，由可信发钥机构完成现实身份核验并签发私钥。通信端不必为每名用户单独查询和
传递证书，业务中的身份规则可以直接映射为密码算法的输入。将格密码与身份基机制结合，可以同时
面向后量子迁移和双方身份约束两项需求。

\subsection{身份基匹配加密的功能意义}

记发送者身份为$s$，发送者指定的接收身份为$r$；接收密钥所属身份为$r'$，接收者认可的发送
身份为$s'$。IB-ME解密成功的必要条件为
\[
  (r=r')\ \land\ (s=s').
\]
表~\ref{tab:identity-cases}列出四种身份组合。只有双方身份同时匹配时，接收端才输出应用消息；
其余三种情况均返回拒绝。

\begin{table}[H]
\centering\small
\caption{接收身份与认可发送身份的匹配关系}
\label{tab:identity-cases}
\begin{tabular}{ccc}
\toprule
\rowcolor{black!8} 接收身份 & 认可的发送身份 & 处理结果 \\
\midrule
匹配 & 匹配 & 恢复消息 \\
匹配 & 不匹配 & 拒绝 \\
不匹配 & 匹配 & 拒绝 \\
不匹配 & 不匹配 & 拒绝 \\
\bottomrule
\end{tabular}
\end{table}

IB-ME与“先签名、后加密”的普通组合相比，重点不在减少一个密码步骤，而在于把两方身份写入统一
的密码关系。接收密钥不能脱离自身身份使用，接收者认可的发送身份也参与密钥派生和签名验证。
若底层加密同时具备身份匿名性，公开信道上的观察者还不能直接从密文中读取通信双方身份。因而，
研究抗量子IB-ME的通用构造与软件实现，既具有密码协议意义，也能够为定向消息系统提供新的实现
方式。

\section{国内外研究现状}

\subsection{身份基加密与分层身份基加密研究}

Shamir于1984年提出身份基密码的基本思想，使公开身份可以直接参与密码运算\cite{shamir1984}。
Boneh和Franklin随后给出基于双线性对的完整身份基加密方案\cite{bonehfranklin2003}。身份基加密
解决了“以身份作为加密目标”的问题，但大型组织往往还需要按照机构、部门、用户或设备层级分配
密钥。Gentry和Silverberg提出分层身份基加密（HIBE），使上层实体能够向下委派身份密钥
\cite{gentrysilverberg2002}。

HIBE的研究逐步从双线性对扩展到格密码。格HIBE通常使用陷门委派，使父级密钥能够生成下一层
身份密钥，同时保持密钥的短向量性质。匿名HIBE进一步要求密文不泄露目标身份路径。LATTE利用
NTRU结构组织量子安全HIBE\cite{zhao2024latte}；Lai等人利用理想格陷门构造分层身份密码
\cite{lai2015trapdoors}。两类方案为本作品实现不同结构的两层匿名HIBE提供了基础。

\subsection{身份基签名研究}

身份基签名（Identity-Based Signature，IBS）允许验证者直接使用签名者身份和系统公开参数
验签。早期研究主要基于双线性对，例如Hess给出的身份基签名方案\cite{hess2003ibs}。格密码研究
发展后，IBS逐步形成两类实现路线：一类通过认证转换把普通格签名接入身份发钥流程；另一类让
身份直接进入模格、环格或矩阵格上的目标关系和短原像采样。

本作品涉及的相关成果包括基于通用认证转换接入ML-DSA和Falcon\cite{galindo2006ibs,
nist2024fips204,falcon2020}，Luo等人的模格IBS\cite{luo2024ibs}，SRGG模格方案
\cite{jia2025ibs}以及Sageloli等人的环格方案\cite{sageloli2023ibs}。这些方案的身份映射、
采样算法、密钥结构和签名尺寸差异明显，为研究通用后端接口和多类格实现提供了有代表性的对象。

\subsection{身份基匹配加密研究}

IB-ME研究如何让发送端和接收端分别指定对方身份。Wang等人给出了具有更强安全目标的IB-ME
构造，并给出由两层匿名HIBE与IBS组成的通用方法及格上实例\cite{wang2024ibme}。在这一方法中，
HIBE的两层身份依次对应接收者和发送者，IBS认证发送者及其生成的协议对象。HIBE与IBS可以分别
选用不同具体方案，因此能够组合多种格结构。

现有研究为IB-ME提供了算法关系和安全模型，但完整软件还要处理底层原语容量不同、对象表示不同、
大型参数资源开销和长期状态管理等问题。本作品以Wang通用构造为理论基础，重点研究多种HIBE与
IBS如何进入同一协议，以及目标参数实现如何完成状态保存、完整消息处理和可重复测量。

\section{现有研究的工程实现问题}

\textbf{原语接口和对象表示不同。} 本作品采用的HIBE与IBS分别使用矩阵、多项式、模向量和NTRU
对象，参数版本、身份映射、采样器和序列化方法也不相同。如果公共协议直接读取后端矩阵或陷门
字段，每增加一种算法都需要修改协议代码，组合数量增加后也难以保证处理逻辑一致。

\textbf{HIBE明文容量与IBS签名尺寸不匹配。} 部分HIBE只适合加密几十字节，而IBS签名可能达到
数KiB、数百KiB甚至更大，应用消息也没有固定长度。将“签名＋消息”直接作为HIBE明文，会使
不少组合超过容量上限。完整系统需要分离短秘密封装和大载荷加密，同时防止两部分被替换或重组。

\textbf{论文参数下的计算和存储开销较高。} 矩阵平方根、离散Gaussian采样和陷门委派若按稠密
矩阵直接实现，可能产生远超常规工作站容量的中间对象。通用JSON表示虽然便于调试，却会重复保存
大量整数的符号、分隔符和十进制字符。实现既要保持原有代数关系，也要减少不必要的矩阵展开和
对象膨胀。

\textbf{单次算法实验不能代表连续消息服务。} 系统建立、用户发钥和固定身份对的叶密钥委派可以
被后续消息复用。如果每条消息都重新执行全部步骤，测试结果会被准备成本主导。软件需要区分长期
状态、会话状态和逐消息对象，并分别测量首次、重复和恢复后的运行成本。

\section{作品研究内容与技术路线}

本作品首先定义HIBE与IBS相互独立的后端接口，并依据两层匿名HIBE与IBS的通用构造实现系统建立、
发送者发钥、接收者发钥、消息加密和匹配解密五个公共算法。协议固定双方身份顺序和认证上下文，
不读取后端内部的矩阵、陷门或采样状态。

其次，作品实现LATTE、Lai两种HIBE，以及BNN＋ML-DSA、BNN＋Falcon、Luo、SRGG、Sageloli
五种常规IBS，形成十种端到端组合。另实现Wang--PW目标参数矩阵SIS后端，并完成其与LATTE的
完整消息流程。这里的Wang--PW是具体IBS后端，不是“扩展IBS”这一类新密码原语。

再次，作品针对实际运行中最突出的三类瓶颈完成Sageloli结构化环采样、SRGG范围感知编码和Lai
接收会话。三项工作分别减少大型中间矩阵、序列化空间和固定身份对的重复委派成本。

最后，作品建设规范对象、状态持久化、命令行和浏览器操作入口，并记录系统建立、双方发钥、首次
与重复收发、错误身份、密文篡改、对象尺寸和峰值内存。图~\ref{fig:work-overview}给出从底层原语
到应用入口的总体技术路线。

\begin{figure}[H]
\centering
\begin{tikzpicture}[node distance=9mm and 10mm]
 \node[flow,text width=34mm] (primitive) {格HIBE与格IBS\\独立后端};
 \node[flow,right=of primitive,text width=34mm] (assembly) {能力声明与兼容性检查\\组合配置};
 \node[flow,right=of assembly,text width=36mm] (protocol) {IB-ME五算法\\混合加密与字节认证};
 \node[check,below=of assembly,text width=34mm] (state) {规范对象\\发钥与状态恢复};
 \node[check,right=of state,text width=36mm] (entry) {命令行与浏览器\\消息收发};
 \draw[flowarrow] (primitive)--(assembly);
 \draw[flowarrow] (assembly)--(protocol);
 \draw[flowarrow] (assembly)--(state);
 \draw[flowarrow] (protocol)--(entry);
 \draw[flowarrow] (state)--(entry);
\end{tikzpicture}
\caption{作品总体技术路线}
\label{fig:work-overview}
\end{figure}

\section{主要工作与创新点}

\begin{table}[H]
\centering\small
\caption{作品的主要工作及其作用}
\label{tab:main-contributions}
\begin{tabularx}{\textwidth}{@{}>{\raggedright\arraybackslash}p{.2\textwidth}>{\raggedright\arraybackslash}p{.37\textwidth}>{\raggedright\arraybackslash}X@{}}
\toprule
\rowcolor{black!8} 工作 & 具体内容 & 解决的问题 \\
\midrule
通用IB-ME设计 & 独立HIBE／IBS接口、五算法、短种子封装、完整字节认证、规范对象 & 使不同格结构和明文容量的原语进入同一协议 \\
多算法组合实现 & 两种HIBE、五种常规IBS的十种配对；Wang--PW目标参数后端与LATTE组合 & 统一实现并比较多类格密码技术路线 \\
关键算法优化 & Sageloli结构化采样、SRGG范围感知编码、Lai接收会话 & 降低矩阵展开、对象存储和重复委派成本 \\
软件与实验系统 & 固定发钥、重启恢复、命令行、浏览器演示、阶段计时、尺寸与内存统计 & 将密码原语转化为可操作、可测量的完整消息系统 \\
\bottomrule
\end{tabularx}
\end{table}

作品的第一项创新是建立面向异构格原语的IB-ME组合框架。新增后端只需实现对应接口、规范编码、
能力声明和自检逻辑，公共五算法与传输格式保持不变。第二项创新是利用环多项式和逐频率低秩分解
实现Sageloli目标参数采样，避免展开巨型稠密矩阵。第三项创新是按照字段实际取值范围编码SRGG
四类对象。第四项创新是为固定身份对建立Lai接收会话，把重复委派移到会话准备阶段。

\section{报告结构}

本报告根据研究内容分为九章，各章安排如下。

\noindent\textbf{第1章 作品概述：}从机构定向通信的双方身份约束出发，说明后量子迁移背景，梳理
HIBE、IBS与IB-ME的研究进展，提出工程实现问题，并概括作品内容与创新点。

\noindent\textbf{第2章 基础知识：}依次介绍格与离散Gaussian、LWE／SIS及其结构变体、身份基加密、
HIBE、IBS和IB-ME，为后续协议和算法实现统一符号与基本概念。

\noindent\textbf{第3章 通用IB-ME系统设计：}给出系统模型、公共对象和五算法伪代码，说明短种子
混合加密、完整密文字节认证、规范编码、正确性和异常处理。

\noindent\textbf{第4章 算法实现与代表方案：}介绍两种HIBE与六种IBS的具体实现、十种常规组合和
Wang--PW目标参数组合，并通过原始指标表与归一化特征图比较A、B、C三组代表及其适用场景。

\noindent\textbf{第5章 关键算法优化：}详细说明Sageloli结构化环采样、SRGG范围感知编码和Lai接收
会话的瓶颈、方法、实现过程与成本模型。

\noindent\textbf{第6章 软件系统与操作流程：}介绍软件分层、稳定接口、新后端接入、固定发钥、状态
恢复、命令行、浏览器入口和部署方式。

\noindent\textbf{第7章 测试与结果分析：}说明实验环境和评价指标，报告完整系统时间、对象尺寸、
通信量、三项优化、功能测试及Wang--PW分阶段结果。

\noindent\textbf{第8章 创新性与应用前景：}集中归纳作品创新，并结合三组代表分析机构消息、受控
终端、固定身份对通信和算法迁移等应用模式及部署条件。

\noindent\textbf{第9章 总结与展望：}总结通用设计、算法实现和实验结论，说明统一同机补测、通信
开销优化和后续后端扩展方向。
